\subsection{Schur's lemma}

In this issue, Hilbert spaces are complex.

PROPOSITION 6 (Schur's Lemma). --- Let \(\varrho\) be a unitary representation of a topological group G in a nonzero Hilbert space E. Then \(\varrho\) is irreducible if and only if \(\mathrm{Hom}_{\mathrm{G}}(\varrho, \varrho)\) is equal to \(\mathbf{C} \cdot 1_{\mathrm{E}}\) .

The space \(\mathrm{Hom}_{\mathrm{G}}(\varrho, \varrho)\) is a unital stellar subalgebra of \(\mathcal{L}(\mathrm{E})\) (lemma 5 of V, p. 380). For every subrepresentation F of E, the orthogonal projector \(p\) with image F is an idempotent element of the stellar algebra \(\mathrm{Hom}_{\mathrm{G}}(\varrho, \varrho)\) (prop. 4 of V, p. 383). If it is equal to \(\mathbf{C} \cdot 1_{\mathrm{E}}\) , this means that \(p = 0\) or \(p = 1_{\mathrm{E}}\) , which means that F = 0 or F = E. Thus \(\varrho\) is irreducible.

Conversely, suppose \(\varrho\) is irreducible. Let \(u\) and \(v\) be commuting elements of the stellar algebra \(\mathrm{Hom}_{\mathrm{G}}(\varrho, \varrho)\) such that \(uv = 0\) . Suppose \(u\) is nonzero. The kernel F of \(u\) then defines a subrepresentation of \(\varrho\) different from E, hence F is reduced to \(\{0\}\) ; since F contains the image of \(v\) , we deduce that \(v = 0\) . By Proposition 10 of I, p. 113, we therefore have \(\mathrm{Hom}_{\mathrm{G}}(\varrho, \varrho) = \mathbf{C} \cdot 1_{\mathrm{E}}\) .

COROLLARY 1. --- Let π be an irreducible unitary representation of a topological group G in a Hilbert space E. Let u be a closed partial operator on E. Suppose that u has dense domain, that dom(u) is stable under π, and that u ∘ π(g) = π(g) ∘ u for all g ∈ G. Then dom(u) = E and u is a scalar multiple of the identity.

The partial operator \(u^{*} \circ u\) is a positive self-adjoint partial operator on E (prop. 12 of IV, p. 241), the same holds for \(v = 1_{\mathrm{E}} + u^{*} \circ u\) and the latter is injective since \(-1 \notin \operatorname{Sp}(u^{*} \circ u)\) (prop. 17 of IV, p. 248). We have \(u^{*} \circ \pi(g) = \pi(g) \circ u^{*}\) for all \(g \in G\) (lemma 6 of V, p. 381), whence \(v \circ \pi(g) = \pi(g) \circ v\) for all \(g \in G\) . Since \(v\) is injective, it follows that \(v^{-1} \circ \pi(g) = \pi(g) \circ v^{-1}\) for all \(g \in G\) . But \(v^{-1}\) belongs to \(\mathcal{L}(\mathrm{E})\) , hence by prop. 6, there exists \(\lambda \in \mathbf{C}\) such that \(v^{-1} = \lambda 1_{\mathrm{E}}\) . Necessarily, \(\lambda \neq 0\) , which implies that \(\mathrm{E} = \mathrm{Im}(v^{-1}) = \mathrm{dom}(v) \subset \mathrm{dom}(u)\) , whence \(\mathrm{dom}(u) = \mathrm{E}\) . Since \(u\) is closed, we have \(u \in \mathcal{L}(\mathrm{E})\) , therefore \(u \in \mathrm{Hom}_{\mathrm{G}}(\pi, \pi)\) and \(u\) is a homothety by Prop. 6.

COROLLARY 2. --- Let \(\varrho\) and \(\pi\) be irreducible unitary representations of a topological group G in Hilbert spaces E and F, respectively. The space \(\mathrm{Hom}_{\mathrm{G}}(\varrho, \pi)\) is of dimension 1 if \(\varrho\) is isomorphic to \(\pi\) and is zero otherwise. In particular, if \(\varrho\) is isomorphic to \(\pi\) , every nonzero G-morphism from \(\varrho\) to \(\pi\) is an isomorphism.

Suppose there exists a nonzero G-morphism \(u\) from \(\varrho\) to \(\pi\) . The linear map \(u^{*} \circ u\) is an element of \(\mathrm{Hom}_{\mathrm{G}}(\varrho, \varrho)\) ; therefore there exists a complex number \(\lambda\) such that \(u^{*} \circ u = \lambda \cdot 1_{\mathrm{E}}\) (prop. 6).

We then have

\[
\langle u (x) \mid u (y) \rangle = \langle x \mid u ^ {*} u (y) \rangle = \lambda \langle x \mid y \rangle
\]

for all \(x\) and \(y\) in E. In particular, \(\lambda \neq 0\) since \(u\) is nonzero. Since \(\| u(x)\| = |\lambda |^{1 / 2}\| x\|\) for all \(x\in \mathrm{E}\) , the linear map \(u\) is injective, and the image of \(u\) is closed in F (Lemma 8 of I, p. 107); it is then a nonzero subrepresentation of the irreducible representation \(\pi\) , whence we deduce that \(u\) is surjective. Thus, \(u\) is an isomorphism from \(\varrho\) to \(\pi\) . The map \(v\mapsto u^{*}\circ v\) is then an isomorphism from the space \(\mathrm{Hom}_{\mathrm{G}}(\varrho ,\pi)\) to the space \(\mathrm{Hom}_{\mathrm{G}}(\varrho ,\varrho)\) , which is of dimension 1 (prop. 6).

COROLLARY 3. --- Let \(\varrho\) be a unitary representation of a topological group G in a Hilbert space E. The representation \(\varrho\) is irreducible if and only if the space \(\mathrm{Sesq}_{\mathrm{G}}(\varrho, \varrho)\) has dimension 1. This space is then generated by the scalar product on E.

In view of prop. 1 of V, p. 381, this follows from prop. 6.

COROLLARY 4. --- Let \(\varrho_{1}\) and \(\varrho_{2}\) be non-isomorphic irreducible unitary representations of a topological group G in Hilbert spaces E \(_{1}\) and E \(_{2}\) . The space Sesq \(_{G}\) ( \(\varrho_{1}, \varrho_{2}\) ) is zero.

By prop. 1 of V, p. 381, this follows from cor. 2.

COROLLARY 5. --- Let \(\varrho\) be a unitary representation of a topological group G in a Hilbert space E. Let \(\pi\) be an irreducible unitary representation of G in a Hilbert space F.

\begin{enumerate}
\def\labelenumi{\alph{enumi})}
\item
  For every nonzero G-morphism u from \(\pi\) to \(\varrho\) , there exists \(\lambda \in \mathbf{R}_{+}^{*}\) such that \(\lambda u\) is isometric;
\item
  Every G-morphism u from π to ρ has closed image;
\item
  Every nonzero G-morphism v from \(\varrho\) to \(\pi\) is surjective.
\end{enumerate}

Let us prove assertion \(a\) ), which implies the assertion \(b\) ) (lemma 8 of I, p. 107). Since the morphism \(u\) is nonzero, it is injective (lemma 2 of V, p. 378). The formula \(q(x,y) = \langle u(x)|u(y)\rangle\) for \(x\) and \(y\) in F then defines a continuous sesquilinear form on F. Since \(u\) is a G-morphism, we have \(q(\pi(g)x,\pi(g)y) = q(x,y)\) for all \(g\) in G and \((x,y)\in\mathrm{F}\times\mathrm{F}\) , hence \(q\) is G-invariant. By Cor. 3, there exists \(\alpha\in\mathbf{R}_{+}^{*}\) such that

\[
q (x, y) = \langle u (x) | u (y) \rangle = \alpha \langle x | y \rangle
\]

for all \((x,y)\in\mathrm{F}\times\mathrm{F}\) , and hence \(\alpha^{-1/2}u\) is isometric.

Let us finally prove assertion \(c\) ). Since \(\pi\) is irreducible, the image of \(v\) is dense in F (Lemma 2 of V, p. 378), hence the adjoint \(v^*\) is injective (TVS, V, p. 41, prop. 4). By \(b\) ), the image H of \(v^*\) is a closed subspace of E; it is therefore a subrepresentation of \(\varrho\) , and \(v^*\), induced by passage to the subspaces, is a G-isomorphism from F to H. The restriction \(w\) of \(v\) to H defines a G-morphism from H to F, which is injective since its kernel is the intersection of H and \(\text{Ker}(v) = \text{H}^\circ\) (loc. cit.). The map \(w\) is therefore an isomorphism (Corollary 2); in particular, \(v\) is surjective.

COROLLARY 6. --- Let \(G_{1}\) and \(G_{2}\) be topological groups. Let \(\varrho_{1}\) be an irreducible unitary representation of \(G_{1}\) in a Hilbert space \(E_{1}\) and \(\varrho_{2}\) an irreducible unitary representation of \(G_{2}\) in a Hilbert space \(E_{2}\) . The external tensor product \(\varrho_{1} \boxtimes \varrho_{2}\) is an irreducible unitary representation of \(G_{1} \times G_{2}\) in \(E_{1} \widehat{\otimes}_{2} E_{2}\) .

The external tensor product \(\varrho = \varrho_{1} \otimes \varrho_{2}\) is a unitary representation of \(G = G_{1} \times G_{2}\) in the space \(E = E_{1} \widehat{\otimes}_{2} E_{2}\) by lemma 8 of V, p. 384.

Let \(q\) be a sesquilinear form \((\mathrm{G}_1 \times \mathrm{G}_2)\)-invariant on E. For every pair \((x_1, y_1) \in \mathrm{E}_1^2\) , the map \((x_2, y_2) \mapsto q(x_1 \otimes x_2, y_1 \otimes y_2)\) belongs to \(\operatorname{Sesq}_{\mathrm{G}_2}(\varrho_2, \varrho_2)\) . Denote by \(b(x_1, y_1)\) the unique complex number such that

\[
q (x _ {1} \otimes x _ {2}, y _ {1} \otimes y _ {2}) = b (x _ {1}, y _ {1}) \left\langle x _ {2} \mid y _ {2} \right\rangle
\]

for all \((x_{2},y_{2})\in\mathrm{E}_{2}^{2}\) (cor. 3). Let \(\varepsilon\in E_{2}\) be of norm 1, so that \(b(x_{1},y_{1})=q(x_{1}\otimes\varepsilon,y_{1}\otimes\varepsilon)\) for all \((x_{1},y_{1})\in\mathrm{E}_{1}^{2}\) . This formula implies that the map b is sesquilinear on \(E_{1}\) . Moreover

\[
\| b (x _ {1}, y _ {1}) \| \leqslant \| q \| \| x _ {1} \otimes \varepsilon \| \| y _ {1} \otimes \varepsilon \| = \| q \| \| x _ {1} \| \| y _ {1} \|
\]

for all \((x_{1},y_{1})\in\mathrm{E}_{1}^{2}\) (EVT, V, p. 26, formula (5)), hence b is continuous.

Let \(g \in G\) and \((x_{1}, y_{1}) \in \mathrm{E}_{1}^{2}\) . Since \(\varrho_{2}\) is unitary and q is invariant, it follows that

\[
b \left(\varrho_ {1} (g) x _ {1}, \varrho_ {1} (g) y _ {1}\right) = q \left(x _ {1} \otimes \varrho_ {2} \left(g ^ {- 1}\right) \varepsilon , y _ {1} \otimes \varrho_ {2} \left(g ^ {- 1}\right) \varepsilon\right) = b \left(x _ {1}, y _ {1}\right),
\]

so that \(b \in \mathrm{Sesq}_{\mathrm{G}_1}(\varrho_1, \varrho_1)\) . There therefore exists a unique \(\lambda \in \mathbf{C}\) such that \(b(x_1, y_1) = \lambda \langle x_1 | y_1 \rangle\) for all \((x_1, y_1) \in \mathrm{E}_1^2\) (cor. 3), that is,

\[
q (x _ {1} \otimes x _ {2}, y _ {1} \otimes y _ {2}) = \lambda \langle x _ {1} | y _ {1} \rangle \langle x _ {2} | y _ {2} \rangle
\]

for all \((x_{1},y_{1},x_{2},y_{2})\in\mathrm{E}_{1}^{2}\times\mathrm{E}_{2}^{2}\) . From this one deduces that the space \(\operatorname{Sesq}_{\mathrm{G}_{1}\times\mathrm{G}_{2}}(\varrho,\varrho)\) has dimension 1, which implies that the representation \(\varrho=\varrho_{1}\boxtimes\varrho_{2}\) is irreducible (loc. cit.).

Recall that U denotes the group of complex numbers of modulus 1.

COROLLARY 7. --- Let \(\pi\) be an irreducible unitary representation of a topological group G in a Hilbert space E. There exists a continuous homomorphism \(\chi\) of the center C of G in U such that \(\pi(z) = \chi(z) \cdot 1_{\mathrm{E}}\) for all \(z \in \mathrm{C}\) . In particular, if G is commutative, we have \(\dim(\mathrm{E}) = 1\) .

For \(z \in \mathrm{C}\) , the map \(\pi(z)\) belongs to \(\mathrm{Hom}_{\mathrm{G}}(\pi, \pi)\) ; it is therefore of the form \(\chi(z) \cdot 1_{\mathrm{E}}\) for some complex number \(\chi(z)\) (prop. 6). Since \(\pi(z)\) is a unitary map, we have \(|\chi(z)| = 1\) . Moreover, since \(\pi\) is a homomorphism, the map \(z \mapsto \chi(z)\) is a homomorphism. Fix \(v \neq 0\) in E; the map \(z \mapsto \chi(z)v = \pi(z)v\) from C to U is continuous, and therefore the homomorphism \(\chi\) is continuous.

Finally, if G is commutative, then C = G, so \(\pi(g) = \chi(g) \cdot 1_{\mathrm{E}}\) for all g in G; every subspace of dimension 1 of E is then a subrepresentation of \(\pi\) , and since \(\pi\) is irreducible, the space E must be of dimension 1.

DEFINITION 6. --- Let π be an irreducible unitary representation of a topological group G in a Hilbert space E. The homomorphism χ of the center C of G into U such that π(z) = χ(z) · \(1_{\mathrm{E}}\) for all z in C is called the central character of π.

Note. --- Let \(\pi\) (resp. \(\varrho\) ) be an irreducible unitary representation of a topological group G (resp. H) in a Hilbert space E (resp. F). Denote by \(\chi\) (resp. \(\eta\) ) the central character of \(\pi\) (resp. \(\varrho\) ). The central character of the representation \(\overline{\pi}\) is \(\overline{\chi}\) and the central character of the irreducible unitary representation \(\pi \boxtimes \varrho\) of G × H (cor. 6) is the character \(\chi \boxtimes \eta\) : \((g, h) \mapsto \chi(g)\eta(h)\) of G × H.

